\chapter{Introduction and conventions}\label{ch:intro}
\chapterstatus{complete}{2026-11-16}{this repository's layout}

Slotty schedules and evaluates student \emph{demos}: short sessions in which a student shows
their work to a teaching assistant (TA) or instructor, who marks it against a rubric. Staff
publish time slots, students book them under a booking policy, staff run the demo day and
record marks, and instructors give marks a final check before students see them.

This document describes Slotty for the people who will audit it, operate it and change it. It is
written to be read in parts: no chapter assumes you have read the ones before it, and every
claim points to the code or configuration it describes.

\section{Who this is for}
\begin{xltabular}{\linewidth}{@{}P{30mm} L L@{}}
  \toprule
  \thead{Reader} & \thead{Needs to} & \thead{Start with}\\
  \midrule
  Auditor (security, architecture, cost) & Verify that controls exist and match intent & \cref{ch:aws,ch:security,ch:decisions,ch:cost}; the \emph{For auditors} boxes\\
  Developer & Change behaviour without breaking rules & \cref{ch:rules,ch:blocks,ch:crosscutting,ch:changes}\\
  Operator & Deploy, watch and repair production & \cref{ch:aws,ch:delivery,ch:ops}\\
  Product owner & Know what the product does for each role & \cref{ch:journeys,ch:roadmap}\\
  AI coding agent & Change Slotty within its rules & \src{CLAUDE.md} and \cref{ch:changes}, then the chapter for the area being changed\\
  \bottomrule
\end{xltabular}

\section{How the document is organised}
The five parts follow the questions an expert asks, in order: \emph{what is it}
(Foundations), \emph{what does it do} (Product), \emph{how is it built} (Software
architecture), \emph{where and how does it run} (Platform and operations), and \emph{how is it
changed and why is it this way} (Lifecycle and governance). The last part also says which chapters
each kind of change must update (\cref{ch:changes}).

\section{Conventions}
\subsection{Callouts}
\begin{brule}{BR-0}
  A rule the system enforces, with where it is enforced.
\end{brule}
\begin{decision}{ADR-000}
  A choice that was made on purpose, with the decision record that explains it.
\end{decision}
\begin{auditnote}
  What to check, and how, to confirm a claim yourself.
\end{auditnote}
\begin{risk}
  A known weakness, with its mitigation or the decision that accepted it.
\end{risk}
\begin{note}
  Context that helps, but changes nothing.
\end{note}
\begin{todo}
  Content still missing. Chapters with these boxes are \emph{draft} or \emph{stub}.
\end{todo}

\subsection{Names in the text}
\begin{xltabular}{\linewidth}{@{}l L@{}}
  \toprule
  \thead{Written as} & \thead{Means}\\
  \midrule
  \src{src/domain/slots.ts} & A file in the repository, linked at the revision this PDF was built from\\
  \code{bookSlot}, \code{FINALIZED} & An identifier in the code: function, model, column or enum value\\
  \route{/courses/[id]} & A URL path in the app\\
  \ui{Open booking} & A label a user sees\\
  \envvar{APP_URL}, \tfvar{web_min} & An environment variable or a Terraform variable\\
  \aws{Amazon RDS} & An AWS service\\
  \bottomrule
\end{xltabular}

\subsection{Diagrams}
All diagrams share one vocabulary (\src{docs/style/slottydiagrams.sty}): teal boxes are Slotty's
own software; boxes with a coloured edge are AWS services, coloured by AWS category (compute
orange, networking purple, database magenta, security and business applications red, storage green,
management pink);
dashed boxes are outside systems. Solid arrows carry requests or data, dashed arrows are queued or
event-driven, and dotted arrows are deployment or management actions.

\section{Building and changing this document}\label{sec:doc-maintenance}
The source is in \src{docs}: \src{docs/slotty.tex} includes one file per chapter from
\src{docs/chapters}, figures are TikZ in \src{docs/figures}, and the design lives in
\src{docs/style}. It builds in a pinned TeX Live container, so every machine and CI produce the
same PDF:
\begin{lstlisting}[language=bash]
bun run docs:build      # → docs/build/slotty.pdf
bun run docs:watch      # rebuild on every save
\end{lstlisting}
CI builds the PDF for every pull request and keeps it as a build artifact.

When you change the system, update the chapters it affects in the same pull request: new or
changed rules in \cref{ch:rules} (with the domain function and its unit test), flows in
\cref{ch:journeys}, infrastructure in \cref{ch:aws}, a new deliberate trade-off as a decision
record in \cref{ch:decisions}. Move a chapter's badge when its status changes, and add a line to
the change log (\cref{app:changelog}).
